\documentclass[11pt,a4paper]{article}
% Preámbulo compartido para los modelos del segundo parcial de ACN.
% Replica el formato del parcial real 2P-2S-2025.
\usepackage[utf8]{inputenc}
\usepackage[T1]{fontenc}
\usepackage[spanish]{babel}
\usepackage{amsmath,amssymb}
\usepackage[margin=2.6cm]{geometry}
\usepackage{enumitem}
\usepackage{array}

\pagestyle{plain}
\setlength{\parindent}{0pt}
\setlength{\parskip}{5pt}

% Listas de incisos: a), b), c), ...
\setlist[enumerate,1]{label=\alph*),leftmargin=2.2em,itemsep=7pt,topsep=7pt}
% Listas numeradas 1., 2., 3., ...
\newlist{puntos}{enumerate}{1}
\setlist[puntos]{label=\arabic*.,leftmargin=2.2em,itemsep=7pt,topsep=7pt}
% Listas de condiciones del enunciado
\setlist[itemize,1]{label=$\bullet$,leftmargin=2em,itemsep=5pt,topsep=6pt}

% Encabezado del parcial
\newcommand{\encabezado}[1]{%
  \begin{center}
    \textbf{\large Aplicaciones Computacionales en Negocios}\\[3pt]
    \textbf{\large Segundo semestre de 2025 --- Segundo parcial}\\[3pt]
    \textbf{\normalsize Modelo de práctica #1}
  \end{center}
  \vspace{4pt}
  \hrule height 0.8pt
  \vspace{14pt}
  Nombre:\par\vspace{12pt}
  Legajo:\par\vspace{12pt}
  DNI:\par\vspace{12pt}
  ¿Te recibís con este parcial?:
  \vspace{10pt}
  \hrule height 0.8pt
  \vspace{16pt}
}

% Título de ejercicio
\newcommand{\ejercicio}[1]{\par\medskip\noindent\textbf{Ejercicio #1.}\ }

% Atajos de notación usados por la cátedra
\newcommand{\R}{\mathbb{R}}
\newcommand{\Z}{\mathbb{Z}}
\newcommand{\dem}{\mathrm{dem}}
\newcommand{\capa}{\mathrm{cap}}
\newcommand{\costo}{\mathrm{costo}}


\begin{document}
\encabezado{3}

%======================================================================
\ejercicio{1}
Consideremos el problema de diseñar el \emph{fixture} de una liga deportiva.
Tenemos un conjunto $E = \{1,\dots,n\}$ de equipos ($n$ es par) y un conjunto
$F = \{1,\dots,2n-2\}$ de fechas. Juegan todos contra todos dos partidos (local
y visitante). Para cada par de equipos $i,j \in E$, $i \neq j$, y cada fecha
$f \in F$, la variable binaria $x_{ijf}$ toma el valor 1 si el equipo $i$ juega
de local contra el equipo $j$ en la fecha $f$. El modelo básico es:

\begin{align}
\sum_{f \in F} x_{ijf} &= 1 && \forall\, i,j \in E,\ i \neq j \\
\sum_{j \in E} \big( x_{ijf} + x_{jif} \big) &= 1 && \forall\, i \in E,\ \forall\, f \in F \\
x_{ijf} &\in \{0,1\} && \forall\, i,j \in E,\ i \neq j,\ \forall\, f \in F \notag
\end{align}

\begin{enumerate}
  \item ¿Qué sucede si eliminamos la restricción (1)? ¿Y si en la restricción
  (2) reemplazamos el signo $=$ por $\le$? Describir en cada caso qué soluciones
  pasan a ser factibles y si el modelo sigue representando el problema.

  \item Agregar una restricción para que el equipo 3 tenga un \emph{patrón
  perfecto}, es decir, para que alterne siempre partidos de local con partidos
  de visitante.

  \item Agregar una restricción para que el equipo 1 juegue al menos cinco
  partidos de local en las primeras ocho fechas, y otra para que el partido
  entre los equipos 2 y 3 (siendo 2 el local) no se juegue entre las fechas 7 y
  12.

  \item Sea $D \subseteq E$ el conjunto de los equipos que pelean por no
  descender de categoría. Agregar una restricción para que en las últimas tres
  fechas se programen al menos cinco partidos entre equipos de este conjunto.
  ¿Por qué puede ser relevante una restricción como ésta?

  \item Supongamos que dividimos $E$ en $n/2$ parejas de equipos, de modo tal
  que los equipos de cada pareja son rivales clásicos. Agregar las restricciones
  necesarias para que haya dos \emph{fechas de clásicos} en las que todas las
  parejas se enfrentan entre sí, sabiendo que las fechas de clásicos no están
  fijadas de antemano: queremos que la primera se planifique entre las fechas 4
  y 12, y la segunda entre las fechas 18 y 30. \emph{Pista:} introducir
  variables binarias asociadas a las fechas.
\end{enumerate}

\newpage
%======================================================================
\ejercicio{2}
Tenemos que conformar un equipo de trabajo a partir de un conjunto $P$ de
candidatos. El equipo debe cubrir un conjunto $R$ de roles, y para cada rol
$r \in R$ se requieren exactamente $D_r \in \Z_+$ personas. Para cada persona
$i \in P$ conocemos su sueldo pretendido $s_i \in \R_+$, y para cada persona
$i \in P$ y cada rol $r \in R$ conocemos el dato $H_{ir} \in \{0,1\}$, que indica
si la persona $i$ está habilitada para desempeñar el rol $r$, y el puntaje
$q_{ir} \in \R_+$ que la persona $i$ obtendría en el rol $r$. Cada persona
seleccionada ocupa exactamente un rol, aunque esté habilitada para más de uno.
El presupuesto salarial total es de $B$ pesos.

\begin{puntos}
  \item Formular un modelo de programación lineal entera para conformar el
  equipo maximizando el puntaje total, respetando las habilitaciones, los
  requerimientos de cada rol y el presupuesto. \emph{Pista:} usar una variable
  binaria $x_{ir}$ por cada persona $i \in P$ y cada rol $r \in R$.

  \item Sea $A \subseteq P$ un conjunto de personas de contratación obligatoria
  y sea $I \subseteq P \times P$ un conjunto de pares de personas que no pueden
  trabajar juntas. Agregar ambas condiciones al modelo.

  \item Decimos que un conjunto $K \subseteq P$ es un \emph{cubrimiento} si
  $\sum_{i \in K} s_i > B$. Explicar por qué, en ese caso, la desigualdad
  \[
  \sum_{i \in K} \sum_{r \in R} x_{ir} \;\le\; |K| - 1
  \]
  es válida para el modelo. ¿Modifica este agregado el conjunto de soluciones
  enteras factibles? ¿Modifica el conjunto de soluciones óptimas? ¿Para qué
  sirve entonces agregarla? ¿Por qué no conviene agregar al modelo todas las
  desigualdades de este tipo?

  \item Supongamos que tres candidatos $u$, $v$ y $w$ son incompatibles de a
  pares, de modo que el modelo del punto 2 incluye las restricciones
  $x_u + x_v \le 1$, $x_v + x_w \le 1$ y $x_u + x_w \le 1$ (escribiendo
  $x_i = \sum_{r \in R} x_{ir}$). Al resolver la relajación lineal, el
  \emph{solver} devuelve $x_u = x_v = x_w = 0.5$. Verificar que esa solución es
  factible para la relajación, proponer una desigualdad válida que la elimine, y
  explicar por qué su agregado acelera la resolución del modelo.
\end{puntos}

\newpage
%======================================================================
\ejercicio{3}
Tenemos un conjunto $T$ de tareas que debemos programar en un conjunto $M$ de
$k > 1$ máquinas. Las máquinas trabajan las 24 horas en forma continua y cada
una procesa una tarea a la vez; una vez que una máquina comienza una tarea, no
se la puede interrumpir hasta que la tarea termine. Es posible ejecutar dos
tareas al mismo tiempo en máquinas distintas, pero no en la misma máquina. Para
cada tarea $i \in T$ y cada máquina $u \in M$ tenemos la duración
$d_{iu} \in \R_+$ de la tarea $i$ si se ejecuta en la máquina $u$, y para cada
tarea $i \in T$ tenemos su fecha de entrega $e_i \in \R_+$, medida a partir del
momento 0. Buscamos determinar en qué máquina y en qué orden programar las
tareas, de modo tal de minimizar la suma de las demoras, entendiendo por demora
de la tarea $i$ el tiempo de terminación de $i$ menos su fecha de entrega,
cuando ese valor es positivo (y cero en caso contrario).

\begin{puntos}
  \item Escribir un modelo de programación lineal entera mixta para este
  problema. \emph{Pista:} usar una variable continua $p_i \ge 0$ por cada tarea
  $i \in T$ que indique en qué momento comienza esa tarea, una variable binaria
  $z_{iu}$ por cada tarea $i \in T$ y cada máquina $u \in M$ que indique si la
  tarea $i$ se asigna a la máquina $u$, y una variable binaria $y_{ij}$ por cada
  par de tareas $i, j \in T$, $i \ne j$, que indique si la tarea $i$ comienza
  antes que la tarea $j$.

  \item Explicar cómo se modela en el objetivo la demora de cada tarea, que es
  de la forma $\max\{0, \cdot\}$. ¿Por qué alcanza con una sola desigualdad por
  tarea? ¿Qué pasaría si el problema fuera de maximización?

  \item Las restricciones de no superposición requieren una constante $M$
  suficientemente grande. Proponer un valor concreto de esa constante en función
  de los datos del problema y explicar por qué conviene que sea lo más chica
  posible, relacionando la respuesta con la relajación lineal del modelo.

  \item Supongamos ahora que tenemos un digrafo $G = (T, A)$ de precedencias
  entre las tareas, de modo tal que si $ij \in A$ entonces la tarea $j$ sólo
  puede comenzar una vez que la tarea $i$ se haya completado. ¿Cómo se modifica
  el modelo? ¿Qué condición debe cumplir $G$ para que el problema sea factible?
\end{puntos}

\end{document}
